\section{The Reversible Action Substrate}
\label{sec:substrate-ra}

A reversible action substrate consists of an action taxonomy, a four-
receipt grammar, an inverse executor for each reversible action, and a
TTL-bounded amendment chain that conditions the active constitution on
the instant the receipt is admitted. The taxonomy and the grammar are
typed; the inverse executors are OS-level operations; the chain is the
substrate-level model that connects the response side to the amendment
side of the prior substrate
\cite{programmableSovereigntyAnonymous}.

\paragraph{Action taxonomy.}
The action taxonomy has twelve variants partitioned into three classes.
Five variants (allow, observe, warn, alert, block) are verdict tags
that do not dispatch to an executor; they are observed by the
admission predicate and recorded by the receipt grammar without OS-
level side effects. Four variants (file quarantine, persistence
disable, process-tree suspend, egress restriction) are reversible:
each has a forward executor and an inverse executor, and the rollback
receipt is admissible when the inverse executor's signed completion
matches the forward executor's recorded effect. Three variants
(process-tree terminate, network isolate, grant revoke) are
destructive: each has a forward executor without an inverse, and
admission reduces to bilateral predicate intersection between a device
polity and an operator polity.

The partition is structural rather than operational. A reversible
variant is one for which an inverse executor exists whose effect, when
applied to the post-effect state, returns the pre-effect state up to
the substrate's notion of equivalence. File rename is the canonical
instance: a forward rename moves a file to a quarantine root, the
inverse rename moves it back to its original path, and the content-
hash gate on the inverse rejects the operation if the original path
has been overwritten in the interim. Process-tree suspend via SIGSTOP
and its inverse SIGCONT is a second instance: the inverse is valid
against the surviving pid set, with the substrate recording which pids
in the original target list have since exited.
Figure~\ref{fig:variants} summarizes the five operational variants and
the shape of the rollback witness each admits.

\begin{figure}[t]
\centering
\small
\begin{tabularx}{\linewidth}{@{}l l X l@{}}
\toprule
\textbf{Variant} & \textbf{Class} & \textbf{Rollback witness} & \textbf{TTL} \\
\midrule
File quarantine        & reversible  & content-hash-gated restore                       & auto-revert \\
Persistence disable    & reversible  & content-hash-gated re-enable                     & auto-revert \\
Process-tree suspend   & reversible  & SIGCONT (substrate-issued capability)            & auto-revert \\
Egress restriction     & reversible  & capability-revoke (substrate-issued)             & auto-revert \\
Process-tree terminate & destructive & absent; bilateral cosignature at admission       & no-revert   \\
\bottomrule
\end{tabularx}
\caption{The five operational action variants and their rollback-witness
shapes. Four reversible variants admit content-hash-gated or
capability-token rollback witnesses constructible at the moment of
admission; the destructive variant reduces admission to bilateral
predicate intersection between a device polity and an operator polity,
since no rollback witness is constructible.}
\label{fig:variants}
\end{figure}

\paragraph{Four-receipt grammar.}
Each action chain emits up to four receipt families. A request receipt
records the action plan: action identifier, action class, target root
node, dry-run flag, TTL duration, rollback reference, reason, issued-
at timestamp, expires-at timestamp, and node and edge counts of the
target subgraph. An execution receipt records the executor's
completion: execution identifier, action identifier, started-at and
completed-at timestamps, an evidence bundle, the actor, and the
effects produced (forward effects only). A rollback receipt closes the
chain when the inverse executor runs: rollback identifier, execution
identifier, action identifier, rolled-back-at timestamp, executor key,
and the inverse effect set with content-hash parity check against the
forward effect set. An acknowledgement receipt records the operator's
attestation that the chain has been observed: acknowledgement
identifier, execution identifier, acknowledged-by, optional note, and
optional control-correlation token carrying the acknowledgement-token
hash.

The receipt grammar is canonical-byte stable. Each receipt body is
serialized under the RFC 8785 canonical JSON discipline the prior
substrate \cite{programmableSovereigntyAnonymous} already requires,
signed under the polity's Ed25519 keypair, and appended to a per-family
ledger. Cross-family references
(rollback's execution identifier, acknowledgement's execution
identifier, control-correlation's response-action identifier) are
checked at verification time by the verifier-owned ledger lookup, not
by trust in the signing kernel.

\paragraph{Bilateral envelopes on destructive variants.}
A destructive action's admission requires a bilateral DSSE envelope
over the canonical request body. The device-polity signature and the
operator-polity signature must be drawn from independent keys; the
strict bilateral DSSE verifier carried by the prior substrate
\cite{programmableSovereigntyAnonymous}
(\codepath{crates/trust/chio-federation/src/bilateral_dsse.rs}) rejects
envelopes with reused signer keys, wrong predicate types, and
noncanonical payloads. The paper-local reduction
\thm{destructive_action_requires_bilateral_admission}
in Section~\ref{sec:model-ra} specializes the proof root's treaty
intersection theorem to the destructive subclass of the action
taxonomy.

\paragraph{TTL-bounded amendment chain.}
A TTL-bounded amendment is a constitutional delta paired with a
positive TTL duration. The pre-amendment constitution is the
substrate's baseline; the post-amendment constitution is the
constitution active inside the TTL window; the auto-reversion at TTL
expiry returns the baseline. The chain is the sequence of such
amendments issued during the substrate's lifetime, indexed by
issued-at timestamp.

The chain is the response-side analog of the prior substrate's
amendment trajectory \cite{programmableSovereigntyAnonymous}. The chain
machinery \thm{essential_preserved_chain} proves that essential
admission is preserved across an amendment sequence whose every step
carries a preservation witness. The TTL-bounded variant adds the auto-reversion:
each step's effect is bounded by the TTL, and the baseline is the
substrate-level invariant the chain returns to whenever the topmost
window expires.

The substrate does not assume real-time semantics. The substrate
treats time as a monotonic integer counter, and the TTL is a duration
in the same counter's units. A real-time scheduler (deployed
separately) advances the counter and triggers expiry callbacks; the
substrate model abstracts the scheduler as an axiom in
Section~\ref{sec:limitations-ra}.
